% GDZ PPN236018647,native576lower–577upper;HKUST577–578.
\NumberedItem{107}
\indent Iisdem diebus circa (Mai. 16.) problema chronologicum de festo paschali eleganter resolvimus.
\par\indent(Promulgatum in \textsc{Zachii} Comm. liter. Aug. 1800, p. 121, 223.)
\par\noindent\hfill[1800 Mai. 16.]

\medskip\begingroup\small
\renewcommand{\thefootnote}{*)}
\indent Der erste, der das Osterfest aus der Jahreszahl mittels arithmetischer Regeln zu bestimmen versucht hat, scheint \textsc{Lambert}\footnote{\textsc{Joh. Heinrich Lambert}, \textit{Einige Anmerkungen über die Kirchenrechnung}, Astronom. Jahrbuch für das Jahr 1778, Berlin 1776, S. 210.} gewesen zu sein, und zwar gilt seine Regel für den julianischen Kalender. Nach \textsc{Carlini}\renewcommand{\thefootnote}{**)}\footnote{\textit{Formole analitiche pel calcolo della Pasqua} di \textsc{Lodovico Ciccolini}, Roma 1817, Biblioteca italiana \textbf{13}, 1819, S. 346. Nach dem anschließenden Briefe ebenda S. 350, \textit{Al signor Francesco Carlini astronomo di Brera} ist \textsc{Carlini} als der Verfasser der Besprechung anzusehen.} sind unter den Vorläufern von \textsc{Gauss} noch \textsc{Oriani}\renewcommand{\thefootnote}{***)}\footnote{\textsc{Barnaba Oriani}, \textit{De usu fractionum continuarum ad inveniendos cyclos calendarii novi et veteris}, Effemeridi di Milano 1786.} und die Patres \textsc{Canovai} und \textsc{del Ricco}\renewcommand{\thefootnote}{\ensuremath{\dagger})}\footnote{\textsc{Stanislao Canovai} e \textsc{Gaëtano del Ricco}, \textit{Elementi di fisica matematica}, Firenze 1788.} zu nennen, die arithmetische Formeln zur Lösung von Aufgaben der Kirchenrechnung aufgestellt haben. Was \textsc{Gauss} im Jahre 1800 von diesen Arbeiten gekannt haben mag, hat sich nicht feststellen
\par\bigskip\noindent\hfill69*\hspace*{.14\linewidth}
\par\endgroup

\clearpage\SourceLeaf{577}\ArithmeticHead{548}{L}{TAGEBUCH MIT ERLÄUTERUNGEN.}
\begingroup\small
\renewcommand{\thefootnote}{*)}
\noindent lassen. Seine allgemeine Osterformel, die sowohl für den julianischen als auch für den gregorianischen Kalender gilt, veröffentlichte er unter dem Titel \textit{Berechnung des Osterfestes} in \textsc{v. Zachs} Monatlicher Correspondenz der Erd- und Himmelskunde, August 1800, S. 121, Verbesserungen S. 223, Werke VI, S. 73; vergl. \textit{Noch etwas über die Bestimmung des Osterfestes}, Braunschweigisches Magazin 12. September 1807, Werke VI, S. 82, ferner \textit{Eine leichte Methode, den Ostersonntag zu finden}, Astronom. Jahrbuch für das Jahr 1814, Berlin 1811, S. 273, und die durch \textsc{P. Tittel} veranlaßte \textit{Berichtigung zu dem Aufsatze: Berechnung des Osterfestes u.s.w.}, Zeitschrift für Astronomie, herausgeg. von \textsc{v. Lindenau} und \textsc{Bohnenberger} \textbf{1}, 1816, S. 158\footnote{Die beiden letztgenannten Aufsätze sind im VI. Bande der Werke nicht enthalten; sie werden mit noch einigen zur Chronologie gehörigen Nachlaßstücken im Bande XI\,1 abgedruckt.}. Daß die \textsc{Gauss}sche Osterregel in ihrer ursprünglichen Form vom Jahre 4200 an ihre Gültigkeit verliert, hat zuerst \textsc{Français}\renewcommand{\thefootnote}{**)}\footnote{\textsc{J. F. Français}, \textit{Solution directe des principaux problèmes du calendrier}, Annales de Mathématiques pures et appliquées \textbf{4}, 1813, S. 273.} öffentlich festgestellt und gleichzeitig die folgende Abänderung angegeben, durch die die \textsc{Gauss}sche Vorschrift unbeschränkte Geltung erlangt: Bezeichnet man den Quotienten, der sich bei der Division der ganzen Zahl \(a\) durch die ganze Zahl \(b\) ergibt, durch \(Q\left(\frac ab\right)\), bedeutet \(A\) die Jahreszahl und setzt man \(k=Q\left(\frac A{100}\right)\), so hat man die Zahl \(p\), die nach \textsc{Gauss} ursprünglicher Festsetzung (Werke VI, S. 78) gleich \(Q\left(\frac k3\right)\) zu nehmen ist, nach \textsc{Français} durch die Formel
\[
p=Q\left(\frac{k-Q\left(\frac{k-17}{25}\right)}3\right)
\]
\noindent zu bestimmen, was mit der von \textsc{Gauss} in der Zeitschrift für Astronomie 1816 angegebenen Festsetzung\renewcommand{\thefootnote}{***)}\footnote{Siehe auch die Werke VI, S. 79, Zeile 5 v. u. abgedruckte handschriftliche Bemerkung.}
\[
p=Q\left(\frac{8k+13}{25}\right)
\]
\noindent übereinstimmt, aber für die Rechnung etwas bequemer ist.
\par\noindent\hfill\textsc{Loewy.}
\par\endgroup
